
\subsection{Masks}
\label{sec:masks}

\begin{figure}[H]
    \centering
    \begin{minipage}[b]{0.32\linewidth}
        \centering
        \includegraphics[width=\linewidth]{masks/mask_quijote_galcut10_cbar.pdf}
    \end{minipage}
    \hfill
    \begin{minipage}[b]{0.32\linewidth}
        \centering
        \includegraphics[width=\linewidth]{masks/mask_quijote_galcut15_cbar.pdf}
    \end{minipage}
    \hfill
    \begin{minipage}[b]{0.32\linewidth}
        \centering
        \includegraphics[width=\linewidth]{masks/mask_quijote_galcut20_cbar.pdf}
    \end{minipage}

    \caption{Analysis masks combining the QUIJOTE low-declination survey mask with 10$^\circ$ (left), 15$^\circ$ (middle), and 20$^\circ$ (right) cuts in galactic latitude.}
    \label{fig:masks}
\end{figure}


\subsection{Foreground model}
\label{sec:foreground_model}


The total angular power spectrum between frequencies $\nu_i$ and $\nu_j$ is modeled as

\begin{equation}
C_\ell(\nu_i, \nu_j) = C_\ell^{\rm sync}(\nu_i, \nu_j) + C_\ell^{\rm dust}(\nu_i, \nu_j) + C_\ell^{\rm corr}(\nu_i, \nu_j),
\end{equation}

where the synchrotron, dust, and synchrotron-dust cross-correlation terms are

\begin{equation}
C_\ell^{\rm sync}(\nu_i, \nu_j) = A_s \left(\frac{\ell}{\ell_{\rm ref}}\right)^{\alpha_s} 
\left(\frac{\nu_i}{\nu_{\rm ref}}\right)^{\beta_s} \left(\frac{\nu_j}{\nu_{\rm ref}}\right)^{\beta_s},
\end{equation}

\begin{equation}
C_\ell^{\rm dust}(\nu_i, \nu_j) = A_d \left(\frac{\ell}{\ell_{\rm ref}}\right)^{\alpha_d} 
S_d(\nu_i) S_d(\nu_j),
\end{equation}

\begin{equation}
\begin{split}
C_\ell^{\rm corr}(\nu_i, \nu_j) = \rho \sqrt{A_s A_d} \Big[ S_s(\nu_i) S_d(\nu_j) + S_s(\nu_j) S_d(\nu_i) \Big] \left(\frac{\ell}{\ell_{\rm ref}}\right)^{(\alpha_s + \alpha_d)/2},
\end{split}
\end{equation}

with $S_s(\nu) = (\nu / \nu_{\rm ref})^{\beta_s}$ and the modified blackbody scaling in Rayleigh-Jeans units

\begin{equation}
S_d(\nu) = \frac{B_\nu(\nu, T_d)}{B_\nu(\nu_0, T_d)} \left(\frac{\nu}{\nu_0}\right)^{\beta_d},
\end{equation}

where $B_\nu$ is the Planck function. Here, $A_s$ and $A_d$ are the amplitudes of the synchrotron and dust components at the reference multipole $\ell_{\rm ref} = 80$ and frequencies $\nu_{\rm ref} = 11.1$~GHz and $\nu_0 = 353.0$~GHz, respectively; $\alpha_s$ and $\alpha_d$ are the multipole spectral indices; $\beta_s$ and $\beta_d$ are the frequency spectral indices; $T_d$ is the dust temperature (fixed at $T_d = 19.6$~K); and $\rho$ is the synchrotron-dust correlation coefficient.

\textcolor{red}{CMB APS}

Before fitting the observed power spectra, we subtract the CMB contribution using the best-fit theoretical $\Lambda$CDM spectrum from Planck 2018 \citep{Planck_2018vi}. This spectrum is obtained from the full Plik TTTEEE likelihood combined with low-multipole (lowE) and gravitational lensing constraints. For polarization modes, we subtract both $C_\ell^{EE}$ and $C_\ell^{BB}$, where the latter includes only the contribution from B-modes generated by gravitational lensing, since primordial B-modes of inflationary origin have not yet been detected. This CMB subtraction is essential to isolate the Galactic foreground emission (synchrotron and thermal dust) that constitutes the target of this analysis.

\subsection{MCMC fitting}
\label{sec:mcmc}


We perform a Markov Chain Monte Carlo (MCMC) fitting procedure using data from QUIJOTE 11 and 17~GHz; WMAP 23, 33, 41, 61, and 94~GHz; and Planck 30, 44, 70, 100, 143, 217, and 353~GHz, covering the multipole range $20 < \ell < 200$. In this study, we use both auto- and cross-spectra between frequency bands. This yields 91 independent spectra from the 14 frequency bands.

\subsection{Bayesian Inference}

We adopt a Bayesian framework to constrain the model parameters $\boldsymbol{\theta} = \{A_s, \alpha_s, \beta_s, A_d, \alpha_d, \beta_d, \rho\}$. The posterior distribution is given by

\begin{equation}
p(\boldsymbol{\theta} \mid \mathbf{D}) \propto \mathcal{L}(\mathbf{D} \mid \boldsymbol{\theta}) \, p(\boldsymbol{\theta}),
\end{equation}

where $\mathcal{L}(\mathbf{D} \mid \boldsymbol{\theta})$ is the likelihood and $p(\boldsymbol{\theta})$ is the prior. The likelihood is constructed from the observed power spectra $\hat{C}_\ell(\nu_i, \nu_j)$ and their uncertainties:

\begin{equation}
\ln \mathcal{L} = -\frac{1}{2} \sum_{i,j,\ell} \frac{[\hat{C}_\ell(\nu_i, \nu_j) - C_\ell(\nu_i, \nu_j; \boldsymbol{\theta})]^2}{\sigma^2_\ell(\nu_i, \nu_j)}.
\end{equation}

For the power-law model fitting, we adopt uniform (flat) priors for all model parameters within physically motivated bounds:

\begin{gather*}
A_s > 0, \\
A_d > 0, \\
\rho \in [-1, 1].
\end{gather*}

For the bin-to-bin analysis, to regularize the fit and incorporate external information from the literature, we impose a Gaussian prior on the synchrotron frequency spectral index:

\begin{equation}
\beta_s \sim \mathcal{N}(-3.1, 0.30^2),
\end{equation}

while maintaining uniform priors for the remaining parameters within the bounds specified above.

\subsection{Fitting methods}
\label{sec:fitting_methods}

We apply three complementary approaches to the angular power spectra computed from the maps:

\paragraph{Power-law fitting}
We repeat the power-law fit using a wider binning, $\Delta\ell = 20$, to test the stability of the fitted parameters with respect to the choice of bin size. Wider bins reduce per-bin noise at the cost of multipole resolution.

\paragraph{Bin-to-bin analysis}
In this approach each multipole bin is fitted independently with $\Delta\ell = 20$, without assuming a global power-law dependence on $\ell$. For each bin we only fit the amplitudes ($A_s$, $A_d$), the frequency spectral indices ($\beta_s$, $\beta_d$), and the correlation ($\rho$); the multipole spectral indices ($\alpha_s$, $\alpha_d$) are not constrained in this analysis. This method allows us to identify scale-dependent behaviour and local deviations from the power-law assumption.


